% GENERATED FILE. Edit canonical lessons, then run npm run generate:books.
% canonical-source-hash: fnv1a64:6c16bd7e32963a5a
% canonical-lessons: LA-C472-colus, LA-C472-fusus, LA-C472-tela, LA-C472-neo, LA-C472-lanificium

\chapter{Distaff, Spindle, Web on the loom, To spin, Wool-working}
\label{ch:la-distaff-spindle-web-on-the-loom-to-spin-wool-working}
\hlchaptermodality{472}

\begin{chapteropening}
\textbf{By the end of this chapter:} \emph{I can talk about spinning and weaving at home: the distaff, the spindle, the web on the loom, spinning, and wool-working.}

Complete the last lesson of chapter 472: I can talk about spinning and weaving at home: the distaff, the spindle, the web on the loom, spinning, and wool-working.
\end{chapteropening}

\section[colus, colī]{colus, colī — a distaff (for spinning)}

\label{lesson:LA-C472-colus}



\begin{quote}
Before the new one: say the Latin for to sharpen, then the Latin for a stonecutter.
\end{quote}

\subsection*{You'll want to know: colus, colī}
\textbf{colus, colī} — \textquotedblleft{}a distaff (for spinning)\textquotedblright{}.

\begin{quote}
Avia colum tenet et lānam trahit.
\end{quote}

\begin{quote}
\emph{Grandmother holds the distaff and draws out the wool.}
\end{quote}

Feminine, with second and fourth declension forms: \emph{colus, colī} or \emph{colūs f.} The stick that holds the raw wool.

\begin{cousinweb}
\emph{colus} is inherited from PIE \emph{k\textsuperscript{w}élus}, built on \emph{k\textsuperscript{w}el-} 'to turn' (EtymDB), the turning stick of the spinner. Lewis \& Short suggest a different root, for being bent or round, so the details are not settled.
\end{cousinweb}

\begin{grammarlens}[title={What you've built}]
One more word for telling, explaining and describing at B1.
\end{grammarlens}

\subsection*{Your turn}
\begin{itemize}
\raggedright
  \item \emph{Say it:} \emph{colus, colī}
  \item \emph{Say it:} \emph{colus, colī}, once more
  \item \emph{Say it:} a sentence of your own with \emph{colus, colī}
  \item \emph{Recall:} say the Latin for to sharpen, then the Latin for a stonecutter, then say \emph{colus, colī} again
\end{itemize}

\begin{tcolorbox}[breakable,colback=teal!4,colframe=teal!35!black,title={Before you move on}]
What does colus, colī mean? (\textquotedblleft{}A distaff (for spinning)\textquotedblright{}.) Say it once more.
\end{tcolorbox}
\section[fūsus, fūsī]{fūsus, fūsī — a spindle}

\label{lesson:LA-C472-fusus}



\begin{quote}
Before the new one: say the Latin for a stonecutter, then the Latin for a distaff.
\end{quote}

\subsection*{You'll want to know: fūsus, fūsī}
\textbf{fūsus, fūsī} — \textquotedblleft{}a spindle\textquotedblright{}.

\begin{quote}
Fūsus in manū ancillae celeriter vertitur.
\end{quote}

\begin{quote}
\emph{The spindle turns fast in the maid's hand.}
\end{quote}

Second declension, masculine: \emph{fūsus, fūsī m.} The weighted stick on which the spun thread is wound.

\begin{cousinweb}
The origin of \emph{fūsus} 'spindle' is not settled: Lewis \& Short give the sense but no source, and EtymDB lists it without a parent. It is a different word from \emph{fūsus} 'poured', from \emph{fundō}.
\end{cousinweb}

\begin{grammarlens}[title={What you've built}]
One more word for telling, explaining and describing at B1.
\end{grammarlens}

\subsection*{Your turn}
\begin{itemize}
\raggedright
  \item \emph{Say it:} \emph{fūsus, fūsī}
  \item \emph{Say it:} \emph{fūsus, fūsī}, once more
  \item \emph{Say it:} a sentence of your own with \emph{fūsus, fūsī}
  \item \emph{Recall:} say the Latin for a stonecutter, then the Latin for a distaff, then say \emph{fūsus, fūsī} again
\end{itemize}

\begin{tcolorbox}[breakable,colback=teal!4,colframe=teal!35!black,title={Before you move on}]
What does fūsus, fūsī mean? (\textquotedblleft{}A spindle\textquotedblright{}.) Say it once more.
\end{tcolorbox}
\section[tēla, tēlae]{tēla, tēlae — a web on the loom, the warp; a loom}

\label{lesson:LA-C472-tela}



\begin{quote}
Before the new one: say the Latin for a distaff, then the Latin for a spindle.
\end{quote}

\subsection*{You'll want to know: tēla, tēlae}
\textbf{tēla, tēlae} — \textquotedblleft{}a web on the loom, the warp; a loom\textquotedblright{}.

\begin{quote}
Fēminae in ātriō ad tēlam labōrant.
\end{quote}

\begin{quote}
\emph{The women work at the loom in the hall.}
\end{quote}

First declension, feminine: \emph{tēla, tēlae f.} Also a spider's web. French \emph{toile} 'cloth' comes from it, and so English \emph{toilet}, first a little cloth.

\begin{cousinweb}
\emph{tēla} is built on \emph{texō} 'to weave' (EtymDB), and Lewis \& Short take it as contracted from an older \emph{texia}, as \emph{āla} stands beside \emph{axilla}. French \emph{toile} 'cloth' comes from it, and English \emph{toilet} from the French diminutive \emph{toilette}.
\end{cousinweb}

\begin{grammarlens}[title={What you've built}]
One more word for telling, explaining and describing at B1.
\end{grammarlens}

\subsection*{Your turn}
\begin{itemize}
\raggedright
  \item \emph{Say it:} \emph{tēla, tēlae}
  \item \emph{Say it:} \emph{tēla, tēlae}, once more
  \item \emph{Say it:} a sentence of your own with \emph{tēla, tēlae}
  \item \emph{Recall:} say the Latin for a distaff, then the Latin for a spindle, then say \emph{tēla, tēlae} again
\end{itemize}

\begin{tcolorbox}[breakable,colback=teal!4,colframe=teal!35!black,title={Before you move on}]
What does tēla, tēlae mean? (\textquotedblleft{}A web on the loom, the warp; a loom\textquotedblright{}.) Say it once more.
\end{tcolorbox}
\section[neō, nēre]{neō, nēre — to spin (thread)}

\label{lesson:LA-C472-neo}



\begin{quote}
Before the new one: say the Latin for a spindle, then the Latin for a web on the loom.
\end{quote}

\subsection*{You'll want to know: neō, nēre}
\textbf{neō, nēre} — \textquotedblleft{}to spin (thread)\textquotedblright{}.

\begin{quote}
Ancillae vesperī lānam nent.
\end{quote}

\begin{quote}
\emph{In the evening the maids spin wool.}
\end{quote}

\emph{neō, nēre, nēvī, nētum}, second conjugation. Short forms: \emph{net} 'she spins', \emph{nent} 'they spin'.

\begin{cousinweb}
\emph{neō} is inherited from Proto-Italic \emph{nēō}, from PIE \emph{(s)neh\textsubscript{1}-} 'to spin' (EtymDB); Lewis \& Short compare Greek \emph{neō}. English \emph{needle} is usually counted a distant Germanic cousin.
\end{cousinweb}

\begin{grammarlens}[title={What you've built}]
One more word for telling, explaining and describing at B1.
\end{grammarlens}

\subsection*{Your turn}
\begin{itemize}
\raggedright
  \item \emph{Say it:} \emph{neō, nēre}
  \item \emph{Say it:} \emph{neō, nēre}, once more
  \item \emph{Say it:} a sentence of your own with \emph{neō, nēre}
  \item \emph{Recall:} say the Latin for a spindle, then the Latin for a web on the loom, then say \emph{neō, nēre} again
\end{itemize}

\begin{tcolorbox}[breakable,colback=teal!4,colframe=teal!35!black,title={Before you move on}]
What does neō, nēre mean? (\textquotedblleft{}To spin (thread)\textquotedblright{}.) Say it once more.
\end{tcolorbox}
\section[lānificium, lānificiī]{lānificium, lānificiī — wool-working, spinning and weaving}

\label{lesson:LA-C472-lanificium}



\begin{quote}
Before the new one: say the Latin for a web on the loom, then the Latin for to spin.
\end{quote}

\subsection*{You'll want to know: lānificium, lānificiī}
\textbf{lānificium, lānificiī} — \textquotedblleft{}wool-working, spinning and weaving\textquotedblright{}.

\begin{quote}
Lānificium in vīllā labor fēminārum est.
\end{quote}

\begin{quote}
\emph{Wool-working in the house is women's work.}
\end{quote}

Second declension, neuter: \emph{lānificium, lānificiī n.}, from \emph{lāna} 'wool' and \emph{faciō} 'to make'. A Roman tomb often praised a woman for it.

\begin{cousinweb}
\emph{lānificium} is built on \emph{lānificus} 'wool-working', from \emph{lāna} 'wool' and \emph{faciō} 'to make' (Lewis \& Short).
\end{cousinweb}

\begin{grammarlens}[title={What you've built}]
One more word for telling, explaining and describing at B1.
\end{grammarlens}

\subsection*{Your turn}
\begin{itemize}
\raggedright
  \item \emph{Say it:} \emph{lānificium, lānificiī}
  \item \emph{Say it:} \emph{lānificium, lānificiī}, once more
  \item \emph{Say it:} a sentence of your own with \emph{lānificium, lānificiī}
  \item \emph{Recall:} say the Latin for a web on the loom, then the Latin for to spin, then say \emph{lānificium, lānificiī} again
\end{itemize}

\begin{tcolorbox}[breakable,colback=teal!4,colframe=teal!35!black,title={Before you move on}]
What does lānificium, lānificiī mean? (\textquotedblleft{}Wool-working, spinning and weaving\textquotedblright{}.) Say it once more.
\end{tcolorbox}
